% losing my mind

\maketitle{}

\tableofcontents{}

\newpage{}

\section{Código fuente}

Se implementaron estructuras de datos abstractas con representaciones concretas distintas, según los requisitos de la tarea. El código fuente está disponible en \url{https://github.com/dontwanttothink/basic_lists}.

\section{Hipótesis}

Se desarrollaron hipótesis sobre la complejidad de tiempo de las operaciones implementadas. En esta sección, tómese $n \in \{0, 1, 2, \dots{}\}$ como la cantidad de elementos en la lista al momento de la operación.

\subsection{Listas enlazadas}

Se implementó una estructura lineal basada en listas enlazadas con los métodos descritos en el documento.

\begin{tcolorbox}[title=Importante]
	Nótese que los métodos \texttt{erase}, \texttt{addBefore} y \texttt{addAfter} se consideraron como respectivamente tomando el nodo a eliminar, el nodo antes del cual agregar un elemento, y el nodo después del cual agregar un elemento. No se toma el valor en sí, lo que permite en ocasiones omitir una búsqueda lineal. (Véase la implementación.)
\end{tcolorbox}

\subsubsection{Lista simplemente enlazada, sin cola}

\begin{table}[htbp]
  \centering{}
  \begin{tabular}{lccc}
    \toprule
    Operación & Peor caso & Caso promedio & Mejor caso \\
    \midrule
    \texttt{pushFront} & $\Theta{} (1)$ & $\Theta{} (1)$ & $\Theta{} (1)$ \\
    \texttt{pushBack} & $\Theta{} (n)$ & $\Theta{} (n)$ & $\Theta{} (n)$ \\
    \texttt{popFront} & $\Theta{} (1)$ & $\Theta{} (1)$ & $\Theta{} (1)$ \\
    \texttt{popBack} & $\Theta{} (n)$ & $\Theta{} (n)$ & $\Theta{} (n)$ \\
    \texttt{find} & $\Theta{} (n)$ & $\Theta{} (n)$ & $\Theta{} (1)$ \\
    \texttt{erase} & $\Theta{} (n)$ & $\Theta{} (n)$ & $\Theta{} (1)$ \\
    \texttt{addBefore} & $\Theta{} (n)$ & $\Theta{} (n)$ & $\Theta{} (1)$ \\
    \texttt{addAfter} & $\Theta{} (1)$ & $\Theta{} (1)$ & $\Theta{} (1)$ \\
    \texttt{isEmpty} & $\Theta{} (1)$ & $\Theta{} (1)$ & $\Theta{} (1)$ \\
    \bottomrule
  \end{tabular}
  \caption{Resumen de las complejidades de tiempo para la lista simplemente enlazada, sin cola}
  \label{tab:resumen-sin-cola}
\end{table}

\begin{description}
	\item[pushFront]{
		Se anticipa $\Theta{} (1)$ en cualquier caso, porque el flujo del algoritmo es idéntico para cualquier $n > 0$ y la cantidad de trabajo realizada es constante.
	}
	\item[pushBack]{
		Se anticipa $\Theta{} (n)$ en cualquier caso, porque para cualquier $n > 0$ se ejecuta una cantidad de trabajo constante por cada elemento de la lista, más una pequeña cantidad de trabajo que es constante e independiente de $n$.
	}
	\item[popFront]{
		Se anticipa $\Theta{} (1)$ en cualquier caso, por razones idénticas a las dadas para \texttt{pushFront}.
	}
	\item[popBack]{
		Se anticipa $\Theta{} (n)$ en cualquier caso, porque para cualquier $n > 0$ se necesita realizar una cantidad constante de trabajo por cada elemento anterior al último, más una cantidad pequeña de tiempo constante e independiente de $n$.
	}
	\item[find]{
		Se anticipa $\Theta{} (n)$ en el peor caso, porque el flujo más costoso ocurre cuando el elemento está al final o no existe, y en aquel caso se realiza una cantidad de trabajo constante por cada elemento de la lista, más una cantidad de trabajo constante independiente de $n$.

		Se anticipa $\Theta{} (1)$ en el mejor caso, porque el flujo menos costoso ocurre cuando el elemento es el primero, y en aquel caso se realiza solo una cantidad de trabajo constante.

		Dado un elemento aleatorio de la lista, se anticipa $\Theta{} (n)$ en el caso promedio, porque la posición esperada del elemento es $\frac{n + 1}{2}$, y se realiza una cantidad de trabajo constante por cada elemento anterior al elemento que se busca, más una cantidad de trabajo constante e independiente de $n$.
	}
	\item[erase]{
		Se anticipa $\Theta{} (n)$ en el peor caso, $\Theta{} (1)$ en el mejor caso, y $\Theta{} (n)$ en el caso promedio, por razones idénticas a las dadas para \texttt{find}, ya que se necesita encontrar el antecesor del nodo a eliminar.
	}
	\item[addBefore]{
		Se anticipa $\Theta{} (n)$ en el peor caso, $\Theta{} (1)$ en el mejor caso, y $\Theta{} (n)$ en el caso promedio, por razones idénticas a las dadas para \texttt{find}.
	}
	\item[addAfter]{
		Se anticipa $\Theta{} (1)$ en cualquier caso, porque para todo $n > 0$ se realiza una cantidad constante de trabajo.
	}
	\item[isEmpty]{
		Se anticipa $\Theta{} (1)$ en cualquier caso, porque la implementación consiste en una verificación simple de igualdad que requiere una cantidad constante de trabajo.
	}
\end{description}

\subsubsection{Lista simplemente enlazada, con cola}

\begin{table}[htbp]
  \centering{}
  \begin{tabular}{lccc}
    \toprule
    Operación & Peor caso & Caso promedio & Mejor caso \\
    \midrule
    \texttt{pushFront} & $\Theta{} (1)$ & $\Theta{} (1)$ & $\Theta{} (1)$ \\
    \texttt{pushBack} & $\Theta{} (1)$ & $\Theta{} (1)$ & $\Theta{} (1)$ \\
    \texttt{popFront} & $\Theta{} (1)$ & $\Theta{} (1)$ & $\Theta{} (1)$ \\
    \texttt{popBack} & $\Theta{} (n)$ & $\Theta{} (n)$ & $\Theta{} (n)$ \\
    \texttt{find} & $\Theta{} (n)$ & $\Theta{} (n)$ & $\Theta{} (1)$ \\
    \texttt{erase} & $\Theta{} (n)$ & $\Theta{} (n)$ & $\Theta{} (1)$ \\
    \texttt{addBefore} & $\Theta{} (n)$ & $\Theta{} (n)$ & $\Theta{} (1)$ \\
    \texttt{addAfter} & $\Theta{} (1)$ & $\Theta{} (1)$ & $\Theta{} (1)$ \\
    \texttt{isEmpty} & $\Theta{} (1)$ & $\Theta{} (1)$ & $\Theta{} (1)$ \\
    \bottomrule
  \end{tabular}
  \caption{Resumen de las complejidades de tiempo para la lista simplemente enlazada, con cola}
  \label{tab:resumen-con-cola}
\end{table}

\begin{description}
	\item[pushFront]{
		Se anticipa $\Theta{} (1)$ en cualquier caso, porque el flujo del algoritmo es idéntico para cualquier $n > 0$ y la cantidad de trabajo realizada es constante.
	}
	\item[pushBack]{
		Se anticipa $\Theta{} (1)$ en cualquier caso por razones idénticas a las dadas para \texttt{pushFront}, gracias al puntero de la cola.
	}
	\item[popFront]{
		Se anticipa $\Theta{} (1)$ en cualquier caso por razones idénticas a las dadas para \texttt{popFront} en la implementación anterior.
	}
	\item[popBack]{
		Se anticipa $\Theta{} (n)$ en todos los casos, por razones idénticas a las dadas para \texttt{popBack} en la implementación anterior.
	}
	\item[find]{
		Se anticipa $\Theta{} (n)$ en el peor caso, $\Theta{} (1)$ en el mejor caso, y $\Theta{} (n)$ en el caso promedio, por razones idénticas a las dadas para \texttt{find} en la implementación anterior.
	}
	\item[erase]{
		Se anticipa $\Theta{} (n)$ en el peor caso, $\Theta{} (1)$ en el mejor caso, y $\Theta{} (n)$ en el caso promedio, por razones idénticas a las dadas para \texttt{erase} en la implementación anterior.
	}
	\item[addBefore]{
		Se anticipa $\Theta{} (n)$ en el peor caso, $\Theta{} (1)$ en el mejor caso, y $\Theta{} (n)$ en el caso promedio, por razones idénticas a las dadas para \texttt{addBefore} en la implementación anterior.
	}
	\item[addAfter]{
		Se anticipa $\Theta{} (1)$ en todos los casos, por razones idénticas a las dadas para \texttt{addAfter} en la implementación anterior.
	}
	\item[isEmpty]{
		Se anticipa $\Theta{} (1)$ en todos los casos, por razones idénticas a las dadas para \texttt{isEmpty} en la implementación anterior.
	}
\end{description}

\subsubsection{Lista doblemente enlazada, sin cola}

\begin{table}[htbp]
  \centering{}
  \begin{tabular}{lccc}
    \toprule
    Operación & Peor caso & Caso promedio & Mejor caso \\
    \midrule
    \texttt{pushFront} & $\Theta{} (1)$ & $\Theta{} (1)$ & $\Theta{} (1)$ \\
    \texttt{pushBack} & $\Theta{} (n)$ & $\Theta{} (n)$ & $\Theta{} (n)$ \\
    \texttt{popFront} & $\Theta{} (1)$ & $\Theta{} (1)$ & $\Theta{} (1)$ \\
    \texttt{popBack} & $\Theta{} (n)$ & $\Theta{} (n)$ & $\Theta{} (n)$ \\
    \texttt{find} & $\Theta{} (n)$ & $\Theta{} (n)$ & $\Theta{} (1)$ \\
    \texttt{erase} & $\Theta{} (1)$ & $\Theta{} (1)$ & $\Theta{} (1)$ \\
    \texttt{addBefore} & $\Theta{} (1)$ & $\Theta{} (1)$ & $\Theta{} (1)$ \\
    \texttt{addAfter} & $\Theta{} (1)$ & $\Theta{} (1)$ & $\Theta{} (1)$ \\
    \texttt{isEmpty} & $\Theta{} (1)$ & $\Theta{} (1)$ & $\Theta{} (1)$ \\
    \bottomrule
  \end{tabular}
  \caption{Resumen de las complejidades de tiempo para la lista doblemente enlazada, sin cola}
  \label{tab:resumen-doble-sin-cola}
\end{table}

\begin{description}
	\item[pushFront]{
		Se anticipa $\Theta{} (1)$ en todos los casos, por razones idénticas a las dadas para \texttt{pushFront} en la implementación anterior.
	}
	\item[pushBack]{
		Se anticipa $\Theta{} (n)$ en todos los casos, porque siempre se ejecuta una cantidad constante de trabajo por cada elemento de la lista, más una pequeña cantidad de trabajo que es constante e independiente de $n$.
	}
	\item[popFront]{
		Se anticipa $\Theta{} (1)$ en todos los casos, por razones idénticas a las dadas para \texttt{pushFront}.
	}
	\item[popBack]{
		Se anticipa $\Theta{} (n)$ en todos los casos, por razones idénticas a las dadas para \texttt{pushBack}.
	}
	\item[find]{
		Se anticipa $\Theta{} (n)$ en el peor caso, $\Theta{} (n)$ en el caso promedio, y $\Theta{} (1)$ en el mejor caso, por razones idénticas a las dadas para \texttt{find} en la implementación anterior.
	}
	\item[erase]{
		Se anticipa $\Theta{} (1)$ en todos los casos, porque se realiza siempre una cantidad constante de trabajo gracias al puntero hacia el elemento anterior.
	}
	\item[addBefore]{
		Se anticipa $\Theta{} (1)$ en todos los casos, por razones idénticas a las dadas para \texttt{erase}.
	}
	\item[addAfter]{
		Se anticipa $\Theta{} (1)$ en todos los casos, por razones idénticas a las dadas para \texttt{addAfter} en la implementación anterior.
	}
	\item[isEmpty]{
		Se anticipa $\Theta{} (1)$ en todos los casos, por razones idénticas a las dadas para \texttt{isEmpty} en la implementación anterior.
	}
\end{description}

\subsubsection{Lista doblemente enlazada, con cola}

\begin{table}[htbp]
  \centering{}
  \begin{tabular}{lccc}
    \toprule
    Operación & Peor caso & Caso promedio & Mejor caso \\
    \midrule
    \texttt{pushFront} & $\Theta{} (1)$ & $\Theta{} (1)$ & $\Theta{} (1)$ \\
    \texttt{pushBack} & $\Theta{} (1)$ & $\Theta{} (1)$ & $\Theta{} (1)$ \\
    \texttt{popFront} & $\Theta{} (1)$ & $\Theta{} (1)$ & $\Theta{} (1)$ \\
    \texttt{popBack} & $\Theta{} (1)$ & $\Theta{} (1)$ & $\Theta{} (1)$ \\
    \texttt{find} & $\Theta{} (n)$ & $\Theta{} (n)$ & $\Theta{} (1)$ \\
    \texttt{erase} & $\Theta{} (1)$ & $\Theta{} (1)$ & $\Theta{} (1)$ \\
    \texttt{addBefore} & $\Theta{} (1)$ & $\Theta{} (1)$ & $\Theta{} (1)$ \\
    \texttt{addAfter} & $\Theta{} (1)$ & $\Theta{} (1)$ & $\Theta{} (1)$ \\
    \texttt{isEmpty} & $\Theta{} (1)$ & $\Theta{} (1)$ & $\Theta{} (1)$ \\
    \bottomrule
  \end{tabular}
  \caption{Resumen de las complejidades de tiempo para la lista doblemente enlazada, con cola}
  \label{tab:resumen-doble-con-cola}
\end{table}

\begin{description}
	\item[pushFront]{
		Se anticipa $\Theta{} (1)$ en cualquier caso, porque el flujo del algoritmo es idéntico para cualquier $n > 0$ y la cantidad de trabajo realizada es constante.
	}
	\item[pushBack]{
		Se anticipa $\Theta{} (1)$ en cualquier caso por razones idénticas a las dadas para \texttt{pushFront}, gracias al puntero de la cola.
	}
	\item[popFront]{
		Se anticipa $\Theta{} (1)$ en cualquier caso por razones idénticas a las dadas para \texttt{pushFront}.
	}
	\item[popBack]{
		Se anticipa $\Theta{} (1)$ en cualquier caso por razones idénticas a las dadas para \texttt{pushBack}.
	}
	\item[find]{
		Se anticipa $\Theta{} (n)$ en el peor caso, $\Theta{} (n)$ en el caso promedio, y $\Theta{} (1)$ en el mejor caso, por razones idénticas a las dadas para \texttt{find} en la implementación anterior.
	}
	\item[erase]{
		Se anticipa $\Theta{} (1)$ en todos los casos, por razones idénticas a las dadas para \texttt{erase} en la implementación anterior.
	}
	\item[addBefore]{
		Se anticipa $\Theta{} (1)$ en todos los casos, por razones idénticas a las dadas para \texttt{addBefore} en la implementación anterior.
	}
	\item[addAfter]{
		Se anticipa $\Theta{} (1)$ en todos los casos, por razones idénticas a las dadas para \texttt{addAfter} en la implementación anterior.
	}
	\item[isEmpty]{
		Se anticipa $\Theta{} (1)$ en todos los casos, por razones idénticas a las dadas para \texttt{isEmpty} en la implementación anterior.
	}
\end{description}

\subsection{Cola dinámica}

Se implementó una cola usando un arreglo. El tamaño inicial del arreglo es 2. Cuando el arreglo se llena, los elementos se copian a un arreglo nuevo con el doble de la capacidad. Cuando el arreglo tiene al menos tres cuartos de su capacidad libre y su capacidad es mayor a 2, los elementos se copian a un arreglo nuevo con la mitad de la capacidad.

A diferencia de las listas enlazadas, tanto esta implementación como la siguiente consideran al método \texttt{delete} como tomando el valor a eliminar en sí. Por lo tanto, ese método involucra una búsqueda lineal.

\begin{table}[htbp]
  \centering{}
  \begin{threeparttable}
    \begin{tabular}{lccc}
      \toprule
      Operación & Peor caso & Mejor caso & Amortizado\tnote{a} \\
      \midrule
      \texttt{enqueue} & $\Theta{} (n)$ & $\Theta{} (1)$ & $\Theta{} (1)$ \\
      \texttt{dequeue} & $\Theta{} (n)$ & $\Theta{} (1)$ & $\Theta{} (1)$ \\
      \texttt{front} & $\Theta{} (1)$ & $\Theta{} (1)$ & ? \\
      \texttt{isEmpty} & $\Theta{} (1)$ & $\Theta{} (1)$ & ? \\
      \texttt{size} & $\Theta{} (1)$ & $\Theta{} (1)$ & ? \\
      \texttt{delete} & $\Theta{} (n)$ & $\Theta{} (1)$ & ? \\
      \bottomrule
    \end{tabular}
    \begin{tablenotes}[flushleft]
      \footnotesize
      \item[a] La complejidad promedio de tiempo de cada operación dada una secuencia de $j$ operaciones. El análisis amortizado se realizó para \texttt{enqueue} y \texttt{dequeue} únicamente.
    \end{tablenotes}
  \end{threeparttable}
  \caption{Resumen de las complejidades de tiempo para la cola dinámica}
  \label{tab:resumen-cola-dinámica}
\end{table}

\begin{description}
	\item[enqueue]{
		Considerando $n \in \{ 2^{a} \mid{} a \in \mathbb{N} \}$, se anticipa $\Theta{} (n)$, porque para $n > 1$ se necesita redimensionar el arreglo, lo cual involucra una cantidad constante de trabajo por cada elemento. Vemos este como el peor caso.

		Considerando $n \notin \{ 2^{a} \mid a \in \mathbb{N}\}$, se anticipa $\Theta{} (1)$, porque se hace una cantidad de trabajo constante. Vemos este como el mejor caso.

		Para determinar la complejidad de tiempo amortizada de $j \in \mathbb{Z}^{+}$ operaciones de \texttt{enqueue}, observamos que realizamos una cantidad de trabajo del orden de $\Theta{}(n)$ en cada operación tal que $n \in \{2^a \mid a \in \mathbb{N}\}$, más una cantidad de trabajo de orden constante en todas las operaciones, necesaria para copiar el elemento nuevo dentro del arreglo. Por lo tanto, para j > 1, la cantidad de tiempo es del orden de
		$$
			j + \sum_{i = 1}^{\lfloor{}\log_{2}(j - 1)\rfloor{}} 2^{i} \text{.}
		$$

		La expresión $\sum_{i = 1}^{\lfloor{}\log_{2}(j - 1)\rfloor{}} 2^{i}$ es una serie geométrica igual a $2^{\lfloor{}\log_{2}(j - 1)\rfloor{} + 1} - 2 = \Theta{}(j)$. Entonces, se anticipa una complejidad promedio de tiempo por operación de $\frac{\Theta{}(j)}{j} = \Theta{}(1)$.
	}
	\item[dequeue]{
		Tómese por $p$ la capacidad del arreglo.

		En el peor caso, se anticipa $\Theta{}(n)$ porque se puede tener que $n - 1 \le \frac{p}{4}$ con $n > 1$; en ese caso, se redimensiona el arreglo, incurriendo en una cantidad de trabajo del orden de $n$.

		Se anticipa $\Theta{}(1)$ en el mejor caso; por ejemplo, si $n - 1 > \frac{p}{4}$, solo se realiza una cantidad constante de trabajo.

		Para una secuencia de $j \in \mathbb{Z}^{+}$ operaciones de \texttt{dequeue}, se anticipa una complejidad promedio de tiempo por operación de $\Theta{}(1)$ a través del razonamiento basado en una función de potencial que se encuentra al final de esta sección.
	}
	\item[front]{
		Se anticipa $\Theta{}(1)$ porque se realiza una cantidad constante de trabajo para todo $n > 0$.
	}
	\item[isEmpty]{
		Se anticipa $\Theta{}(1)$ porque se realiza una cantidad constante de trabajo invariablemente.
	}
	\item[size]{
		Se anticipa $\Theta{}(1)$ porque se realiza una cantidad constante de trabajo invariablemente.
	}
	\item[delete]{
		Se anticipa $\Theta{}(n)$ en el peor caso porque los flujos más costosos posibles requieren una cantidad de trabajo del orden de $n$. La implementación busca al elemento desde el principio; entonces, si el elemento a eliminar es el último, se realizará una cantidad constante de trabajo por cada elemento de la estructura. Si $n - 1$ es a lo más un cuarto de la capacidad del arreglo y $n > 1$, el arreglo es redimensionado, realizando otra vez una cantidad de trabajo del orden de $n$. Si se elimina un elemento en el medio de la cola, se mueven todos los elementos a la izquierda del elemento eliminado, lo que incurre también en una cantidad de trabajo del orden de $n$.

		Un mejor caso ocurre cuando el elemento a eliminar es el primero y además $n - 1$ es mayor que un cuarto de la capacidad del arreglo. Se anticipa $\Theta{}(1)$ en ese caso, porque solo se realiza una cantidad de trabajo constante.
	}
\end{description}

Se anticipa que una secuencia arbitraria de longitud $j > 0$ compuesta de \texttt{enqueue} y \texttt{dequeue} tenga una complejidad promedio de tiempo por operación de $\Theta{} (1)$, con la siguiente justificación:

Sean $\left[D_i\right]_{i \in \{ 0, 1, 2, \dots, j\}}$ los estados de la estructura \emph{después} de $i$ operaciones. Usando el método de potenciales, definimos
$$
	\Phi{}(D_i) = \begin{cases}
		2\left(n - \frac{p}{2}\right) & \text{si } n > \frac{p}{2} \\
		\frac{p}{2} - n & \text{si } 0 < n \le \frac{p}{2} \\
		0 & \text{si } n = 0 \\
	\end{cases}
$$
donde $n$ es el número de elementos en $D_i$ y $p$ es su capacidad.

Entonces, se tiene que $\Phi{}(D_0) = 0$ porque la cola inicia sin elementos. Además, $\Phi{}(D_i)$ es no negativo para cualquier $i \in \{ 0, 1, 2, \dots, j\}$.

Ahora consideramos el costo amortizado de cada operación en la secuencia. Otra vez, tomamos $n$ como el número de elementos en la cola $D_{i}$. Si la $i$-ésima operación es del tipo \texttt{enqueue}, existen dos casos: puede haber redimensionamiento o no. Si el arreglo se redimensiona, tomamos el costo real como $n$. Por lo tanto, el costo amortizado está dado por
\begin{align*}
	n + \Phi{}(D_{i}) &- \Phi{}(D_{i - 1}) \\
	&= n + 2\left(n - \frac{2(n - 1)}{2}\right) - 2\left(n - 1 - \frac{n - 1}{2}\right) \\
	&= n + \left(2\right) - \left(n - 1\right) \\
	&= 3 = \Theta{}(1)
\end{align*}
ya que $p = n - 1$ inmediatamente antes de redimensionar y $p = 2(n - 1)$ después de redimensionar. Si el arreglo no se redimensiona, tomamos el costo real como 1 y notamos que el potencial aumenta a lo más dos unidades.\footnote{A partir de la definición de la función de potencial, se puede ver que el potencial disminuye una unidad si se ejecuta \texttt{enqueue} y resulta $1 < n \le \frac{p}{2}$, aumenta dos unidades si resulta $n > \frac{p}{2}$ y no cambia si $n = 1$.} Entonces, se tiene $1 + \Phi{}(D_{i}) - \Phi{}(D_{i - 1}) \le 1 + 2 = 3$, por lo que $1 + \Phi{}(D_{i}) - \Phi{}(D_{i - 1}) = O(1)$ en este caso también.\footnote{Formalmente, no se tiene que la complejidad amortizada de \texttt{enqueue} sea $\Theta{}(1)$ porque no existen $k \in \mathbb{R}^{+}$ e $I \in \mathbb{Z}^{+}$ tales que $1k < 1 + \Phi{}(D_{i}) - \Phi{}(D_{i - 1})$ para todo $i > I$, puesto que el costo amortizado puede no ser positivo. Sin embargo, $O(1)$ es suficiente para el argumento. Esta nota aplica también a \texttt{dequeue}.}

Paralelamente, si la $i$-ésima operación es del tipo \texttt{dequeue}, consideramos los mismos dos casos. Si el arreglo se redimensiona, tomamos otra vez el costo real como $n$. Por lo tanto, el costo amortizado está dado por
\begin{align*}
	n + \Phi{}(D_{i}) &- \Phi{}(D_{i - 1}) \\
	&= n + \left(\frac{2n}{2} - n\right) - \left(\frac{4n}{2} - (n + 1)\right) \\
	&= n + \left(0\right) - \left(n - 1\right) \\
	&= 1 = \Theta{}(1)
\end{align*}
ya que $p = 4n$ inmediatamente antes de redimensionar y $p = 2n$ después de redimensionar. Si no hay redimensionamiento, entonces tomamos el costo real como $1$ y notamos que el potencial aumenta a lo más una unidad.\footnote{Esta vez, vemos que el potencial disminuye dos unidades si resulta $n \ge \frac{p}{2}$, aumenta una unidad si resulta $n < \frac{p}{2}$ y no cambia si resulta $n = 0$.} Entonces, se tiene $1 + \Phi{}(D_{i}) - \Phi{}(D_{i - 1}) \le 1 + 1 = 2$, por lo que $1 + \Phi{}(D_{i}) - \Phi{}(D_{i - 1}) = O(1)$ en este caso también.

El costo amortizado de ambas operaciones es $O(1)$ y por lo tanto el costo total amortizado de $j$ operaciones es $O(j)$. Dado que $\Phi{}(D_i) \ge \Phi{}(D_0)$ para todo $i \in \{ 0, 1, 2, \dots{}, j\}$, el costo amortizado total de $j$ operaciones es igual o mayor al costo total real. Se deduce que el costo real en el peor caso de $j$ operaciones es entonces $O(j)$, y dado que cada operación tiene un costo real entero distinto de cero, se deduce también que el costo de $j$ operaciones es $\Omega{}(j)$. Por lo tanto, se anticipa que el costo de $j$ operaciones sea $\Theta{}(j)$ y, equivalentemente, que el costo promedio de cada operación en una secuencia de $j$ operaciones sea $\frac{\Theta{}(j)}{j} = \Theta{}(1)$.

\subsection{Pila dinámica}
\label{sec:stack}

La pila se implementó con un arreglo usando la misma estrategia de redimensionamiento que la cola dinámica.

\begin{table}[htbp]
  \centering{}
  \begin{threeparttable}
    \begin{tabular}{lccc}
      \toprule
      Operación & Peor caso & Mejor caso & Amortizado\tnote{a} \\
      \midrule
      \texttt{push} & $\Theta{} (n)$ & $\Theta{} (1)$ & $\Theta{} (1)$ \\
      \texttt{pop} & $\Theta{} (n)$ & $\Theta{} (1)$ & $\Theta{} (1)$ \\
      \texttt{peek} & $\Theta{} (1)$ & $\Theta{} (1)$ & ? \\
      \texttt{isEmpty} & $\Theta{} (1)$ & $\Theta{} (1)$ & ? \\
      \texttt{size} & $\Theta{} (1)$ & $\Theta{} (1)$ & ? \\
      \texttt{delete} & $\Theta{} (n)$ & $\Theta{} (n)$ & ? \\
      \bottomrule
    \end{tabular}
    \begin{tablenotes}[flushleft]
      \footnotesize
      \item[a] La complejidad promedio de tiempo de cada operación dada una secuencia de $j$ operaciones. El análisis amortizado se realizó para \texttt{push} y \texttt{pop} únicamente.
    \end{tablenotes}
  \end{threeparttable}
  \caption{Resumen de las complejidades de tiempo para la pila dinámica}
  \label{tab:resumen-pila-dinámica}
\end{table}

\begin{description}
	\item[push]{
		Por razones idénticas a las dadas para \texttt{enqueue} en la cola dinámica, se anticipa $\Theta{}(1)$ en el mejor caso, $\Theta{}(n)$ en el peor caso, y una complejidad de tiempo promedio de $\Theta{}(1)$ por operación dada una secuencia de $j \in \mathbb{Z}^{+}$ operaciones.
	}
	\item[pop]{
		Se necesita una cantidad constante de trabajo excepto cuando se realiza un redimensionamiento. Por razones idénticas a las dadas para \texttt{dequeue} en la cola dinámica, se anticipa $\Theta{}(1)$ en el mejor caso y $\Theta{}(n)$ en el peor caso. Se anticipa una complejidad de tiempo promedio de $\Theta{}(1)$ por operación dada una secuencia de $j \in \mathbb{Z}^{+}$ operaciones, porque el análisis amortizado para la cola dinámica es aplicable también a la pila dinámica.
	}
	\item[peek]{
		Para cualquier $n > 0$, se necesita una cantidad de trabajo constante. Por lo tanto, se anticipa $\Theta{}(1)$.
	}
	\item[isEmpty]{
		Se necesita siempre una cantidad de trabajo constante; entonces, se anticipa $\Theta{}(1)$.
	}
	\item[size]{
		Se anticipa $\Theta{}(1)$ por razones idénticas a las dadas para \texttt{isEmpty}.
	}
	\item[delete]{
		Se anticipa $\Theta{}(n)$ en todos los casos, porque la búsqueda se debe hacer desde el principio y luego se necesitan mover todos los elementos que estén a la derecha del elemento eliminado. La búsqueda necesita una cantidad de trabajo constante por cada elemento a la izquierda del elemento a eliminar, y la traslación necesita una cantidad de trabajo constante por cada elemento restante a la derecha. También es posible que se necesite hacer un redimensionamiento, contribuyendo otra vez una cantidad de trabajo del orden de $n$.
	}
\end{description}

El análisis amortizado dado para secuencias de operaciones de \texttt{enqueue} y \texttt{dequeue} en la cola dinámica es aplicable también a secuencias de operaciones de \texttt{push} y \texttt{pop} porque la estrategia de redimensionamiento es idéntica. Por lo tanto, se anticipa que una secuencia cualquiera de $j \in \mathbb{Z}^{+}$ operaciones de \texttt{push} y \texttt{pop} tenga una complejidad de tiempo de $\Theta{}(j)$ y por lo tanto la complejidad de tiempo promedio por operación sea $\Theta{}(1)$.

\section{Mediciones de tiempo}

Se ejecutaron los métodos de cada implementación para ciertos $n$ desde $1$ hasta $2^{26} = 67\,108\,864$. En general, se tomaron mediciones de los métodos en cada poder de dos.

Se tomaron mediciones adicionales para los métodos \texttt{push} y \texttt{enqueue} de las estructuras basadas en arreglos dinámicos para permitir observar el mejor caso. Los puntos de medición de estos métodos también fueron exponenciales, pero ocho veces más frecuentes. Es decir, se realizaron mediciones cada poder de $2^{\frac{1}{8}}$.

Los métodos que reciben argumentos, como \texttt{find}, se ejecutaron con una entrada elegida aleatoriamente dentro de los elementos existentes de la estructura.

Se repitió todo el proceso 20 veces y se tomaron solo los valores medios.

Las mediciones se realizaron en nanosegundos, usando \texttt{System.nanoTime()}, porque era el mayor nivel de precisión disponible. Medir la cantidad de tiempo en milisegundos no era adecuado porque la mayoría de las mediciones tardó significativamente menos de un milisegundo.

\subsection{Listas enlazadas}

\subsubsection{Lista simplemente enlazada, sin cola}

\begin{figure}[htbp]
  \centering
  \includegraphics[width=\linewidth]{list_sll_overlay.pdf}
  \caption{Tiempos de ejecución superpuestos de la lista simplemente enlazada, sin cola}
  \label{fig:list_sll_overlay}
\end{figure}

La figura~\ref{fig:list_sll_overlay} permite observar el crecimiento de los distintos tiempos de ejecución a medida que $n$ aumenta. Todos los resultados son consistentes con lo anticipado, excluyendo algo de ruido al inicio y al final.

Los métodos \texttt{isEmpty}, \texttt{popFront}, \texttt{pushFront} y \texttt{addAfter} no muestran una correlación con la cantidad de elementos en la estructura, lo cual obedece a la expectativa de $\Theta{}(1)$.

Los demás métodos muestran una relación lineal clara, como también se había predecido. En particular, el rendimiento lineal de \texttt{addBefore} y \texttt{find} es esperado porque el nodo y el elemento, respectivamente, fueron escogidos aleatoriamente dentro de los elementos existentes de la lista.

\subsubsection{Lista simplemente enlazada, con cola}

\begin{figure}[htbp]
  \centering
  \includegraphics[width=\linewidth]{list_sllt_overlay.pdf}
  \caption{Tiempos de ejecución superpuestos de la lista simplemente enlazada, con cola}
  \label{fig:list_sllt_overlay}
\end{figure}

La figura~\ref{fig:list_sllt_overlay} muestra los cambios que se esperaban al agregar un puntero hacia la cola.

La única diferencia significativa es el rendimiento de \texttt{pushBack}, que ahora parece $\Theta{}(1)$, como era esperado. El puntero hacia la cola permite agregar un elemento al final de la lista sin necesidad de recorrerla toda.

El resto de las operaciones también coinciden. Por ejemplo, otra vez vemos el crecimiento lineal de \texttt{popBack}; la línea es casi recta.

\subsubsection{Lista doblemente enlazada, sin cola}

\begin{figure}[htbp]
  \centering
  \includegraphics[width=\linewidth]{list_dll_overlay.pdf}
  \caption{Tiempos de ejecución superpuestos de la lista doblemente enlazada, sin cola}
  \label{fig:list_dll_overlay}
\end{figure}

La figura~\ref{fig:list_dll_overlay} muestra los comportamientos claramente lineales de \texttt{find}, \texttt{pushBack} y \texttt{popBack} que se anticipaban. Estos métodos necesitan realizar una iteración del orden de $n$.

Al mismo tiempo, el resto de los métodos muestran una línea casi recta. Por ejemplo, \texttt{erase} es ahora de tiempo constante, gracias al puntero hacia el elemento anterior.

Las mediciones coinciden con todas las expectativas.

\subsubsection{Lista doblemente enlazada, con cola}

\begin{figure}[htbp]
  \centering
  \includegraphics[width=\linewidth]{list_dllt_overlay.pdf}
  \caption{Tiempos de ejecución superpuestos de la lista doblemente enlazada, con cola}
  \label{fig:list_dllt_overlay}
\end{figure}

La figura~\ref{fig:list_dllt_overlay} muestra un comportamiento lineal solo para \texttt{find}, mientras que el resto de las operaciones se mantienen más o menos iguales a medida que $n$ aumenta. Esto es idéntico a lo sugerido por el cuadro~\ref{tab:resumen-doble-con-cola}. La combinación de nodos doblemente enlazados y un puntero hacia la cola permite realizar casi todas las operaciones en tiempo constante, siempre que no sea necesario hacer una búsqueda primero.

\subsection{Cola dinámica}

\begin{figure}[htbp]
  \centering
  \includegraphics[width=\linewidth]{array_queue_overlay.pdf}
  \caption{Tiempos de ejecución superpuestos de la cola dinámica}
  \label{fig:array_queue_overlay}
\end{figure}

En la figura~\ref{fig:array_queue_overlay} se aprecia tanto el comportamiento lineal como el comportamiento constante del método \texttt{enqueue}. En los poderes de dos, el arreglo tiene que copiarse, incurriendo en una cantidad de trabajo del orden de $n$. Las aspas azules muestran este fenómeno. Por el contrario, el resto de las mediciones para \texttt{push} muestran una cantidad de tiempo acotada y relativamente constante.

Ya que los puntos medidos no coincidían con los puntos de redimensionamiento para \texttt{dequeue}, solo se ven los tiempos de ejecución de su mejor caso.

Como era esperado, \texttt{find} continúa tomando una cantidad de tiempo del orden de $n$, por las mismas razones que en las otras estructuras.

El resto de las operaciones parece tomar una cantidad de tiempo que no aumenta a medida que aumenta $n$, que es lo que se anticipaba.

\subsection{Pila dinámica}

\begin{figure}[htbp]
  \centering
  \includegraphics[width=\linewidth]{array_stack_overlay.pdf}
  \caption{Tiempos de ejecución superpuestos del arreglo dinámico (pila)}
  \label{fig:array_stack_overlay}
\end{figure}

La pila dinámica es muy parecida a la cola dinámica. En la figura~\ref{fig:array_stack_overlay} se aprecia tanto el comportamiento lineal como el comportamiento constante del método \texttt{push}. En los poderes de dos, el arreglo tiene que copiarse, incurriendo en una cantidad de trabajo del orden de $n$. Las aspas azules muestran este fenómeno. Se puede ver que el resto de las mediciones para \texttt{push} no aumentan con $n$.

Ya que los puntos medidos no coincidían con los puntos de redimensionamiento para \texttt{pop}, solo se ven los tiempos de ejecución de su mejor caso.

Una vez más, \texttt{find} continúa tomando una cantidad de tiempo del orden de $n$.

El resto de las operaciones parece tomar una cantidad de tiempo que no aumenta a medida que aumenta $n$, que es lo que se anticipaba.

\subsection{Comparación entre los métodos análogos de las listas enlazadas y las colas y arreglos dinámicos}

\begin{figure}[htbp]
  \centering
  \includegraphics[width=\linewidth]{mystack_comparativa_push.pdf}
  \caption{Comparación de los tiempos de ejecución de \texttt{push} (pila dinámica) y \texttt{pushBack} (lista enlazada simplemente con cola)}
  \label{fig:mystack_comparativa_push}
\end{figure}

La figura~\ref{fig:mystack_comparativa_push} muestra que el comportamiento del método \texttt{pushBack} de la lista simplemente enlazada con cola es similar al del método \texttt{push} de la pila dinámica, salvo cuando la pila dinámica necesita redimensionarse. Sin embargo, el método de la pila dinámica tarda significativamente menos en el mejor caso que el método de la lista enlazada.

Se eligió la lista simplemente enlazada con cola para esta comparación porque ofrece una complejidad de tiempo de $\Theta{}(1)$ para \texttt{pushBack}, igual a la complejidad de tiempo promedio por operación de la operación \texttt{push} de la pila dinámica.

\begin{figure}[htbp]
  \centering
  \includegraphics[width=\linewidth]{mystack_comparativa_pop.pdf}
  \caption{Comparación de los tiempos de ejecución de \texttt{pop} (pila dinámica) y \texttt{popBack} (lista enlazada doblemente con cola)}
  \label{fig:mystack_comparativa_pop}
\end{figure}

La figura~\ref{fig:mystack_comparativa_pop} muestra que el comportamiento a medida que $n$ aumenta del método \texttt{pop} de la pila dinámica y del método \texttt{popBack} de la lista doblemente enlazada con cola es casi igual, al menos para los valores de $n$ medidos. Esta gráfica no incluye el peor caso para \texttt{pop} porque el redimensionamiento no es necesario excepto en las condiciones específicas descritas en la sección~\ref{sec:stack}, y estas condiciones no se dieron durante las pruebas.

Se eligió la lista doblemente enlazada con cola porque ofrece una complejidad de tiempo de $\Theta{}(1)$ para esta operación, lo cual iguala la complejidad de tiempo en el mejor caso de la operación \texttt{pop} de la pila dinámica.

\begin{figure}[htbp]
  \centering
  \includegraphics[width=\linewidth]{mystack_comparativa_delete.pdf}
  \caption{Comparación de los tiempos de ejecución de \texttt{delete} (pila dinámica) y \texttt{find} + \texttt{erase} (lista enlazada simplemente sin cola)}
  \label{fig:mystack_comparativa_delete}
\end{figure}

La figura~\ref{fig:mystack_comparativa_delete} permite comparar el rendimiento de la pila dinámica y la lista enlazada simplemente sin cola al eliminar un elemento aleatorio. Debido a que la lista enlazada necesita una referencia al nodo a eliminar, eliminar la primera ocurrencia de un \emph{elemento} dado requiere realizar una búsqueda lineal primero. Para aproximar el tiempo total de ambas operaciones, se sumaron las medidas del método \texttt{find} y el método \texttt{erase}.

El resultado es consistente con el comportamiento esperado de $\Theta{}(n)$ para ambas rutinas.

Se eligió una lista simplemente enlazada sin cola porque todas las implementaciones de listas tardan una cantidad de tiempo del orden de $n$ en buscar un elemento dado.

\begin{figure}[htbp]
  \centering
  \includegraphics[width=\linewidth]{myqueue_comparativa_enqueue.pdf}
  \caption{Comparación de los tiempos de ejecución de \texttt{enqueue} (cola dinámica) y \texttt{pushBack} (lista enlazada simplemente con cola)}
  \label{fig:myqueue_comparativa_enqueue}
\end{figure}

Parecido al caso de la pila dinámica, la figura~\ref{fig:myqueue_comparativa_enqueue} muestra que el comportamiento del método \texttt{pushBack} de la lista simplemente enlazada con cola es similar al del método \texttt{enqueue} de la cola dinámica, salvo cuando la cola dinámica necesita redimensionarse. Sin embargo, el método de la pila dinámica tarda significativamente menos en el mejor caso que el método de la lista enlazada.

Se eligió la lista simplemente enlazada con cola porque ofrece una complejidad de tiempo de $\Theta{}(1)$ para \texttt{pushBack}, igual a la complejidad de tiempo promedio por operación de la operación \texttt{enqueue} de la cola dinámica.

\begin{figure}[htbp]
  \centering
  \includegraphics[width=\linewidth]{myqueue_comparativa_dequeue.pdf}
  \caption{Comparación de los tiempos de ejecución de \texttt{dequeue} (cola dinámica) y \texttt{popFront} (lista enlazada simplemente sin cola)}
  \label{fig:myqueue_comparativa_dequeue}
\end{figure}

La figura~\ref{fig:myqueue_comparativa_dequeue} muestra que el comportamiento a medida que $n$ aumenta del método \texttt{dequeue} de la pila dinámica y del método \texttt{popFront} de la lista doblemente enlazada con cola es casi igual, al menos para los valores de $n$ medidos. Esta gráfica no incluye el peor caso para \texttt{dequeue} porque el redimensionamiento no es necesario excepto en las mismas condiciones específicas discutidas para la pila dinámica.

Se eligió la lista doblemente enlazada con cola porque ofrece una complejidad de tiempo de $\Theta{}(1)$ para esta operación, lo cual iguala la complejidad de tiempo en el mejor caso de la operación \texttt{pop} de la cola dinámica.

\begin{figure}[htbp]
  \centering
  \includegraphics[width=\linewidth]{myqueue_comparativa_delete.pdf}
  \caption{Comparación de los tiempos de ejecución de \texttt{delete} (cola dinámica) y \texttt{find} + \texttt{erase} (lista enlazada simplemente sin cola)}
  \label{fig:myqueue_comparativa_delete}
\end{figure}

La figura~\ref{fig:myqueue_comparativa_delete} permite comparar el rendimiento de la cola dinámica y la lista simplemente enlazada sin cola al eliminar un elemento aleatorio. De nuevo, debido a que la lista enlazada necesita una referencia al nodo a eliminar, eliminar la primera ocurrencia de un \emph{elemento} dado requiere realizar una búsqueda lineal primero. Para aproximar el tiempo total de ambas operaciones, se sumaron las medidas del método \texttt{find} y el método \texttt{erase}.

El resultado es consistente con el comportamiento esperado de $\Theta{}(n)$ para ambas rutinas.

Se eligió una lista simplemente enlazada sin cola porque todas las implementaciones de listas tardan una cantidad de tiempo del orden de $n$ en buscar un elemento dado.

\section{Conclusiones}

En cuanto a las complejidades de tiempo promedio por operación, las listas enlazadas superan a los arreglos dinámicos solamente cuando se necesitan hacer inserciones en puntos arbitrarios de la lista. Cuando solo se realizan inserciones al inicio o al final, una lista doblemente enlazada con cola es superior a un arreglo dinámico únicamente cuando las pausas de redimensionamiento no son aceptables.

Se escucha a menudo que las listas enlazadas no son apropiadas en la gran mayoría de las situaciones prácticas. Por ejemplo, la \href{https://doc.rust-lang.org/std/collections/struct.LinkedList.html}{librería estandar del lenguaje de programación Rust} advierte que ``es casi siempre mejor usar un Vec [pila dinámica] o un VecQueue [cola dinámica] porque los contenedores basados en arreglos son generalmente más rápidos, más eficientes en memoria, y hacen un mejor uso del caché del procesador.''

Una cola dinámica puede ser útil en la practicidad para eficientemente modelar la lista de pedidos pendientes de un restaurante, porque las órdenes realizadas primero suelen empezar a prepararse antes que las órdenes realizadas después.

Una lista doblemente enlazada puede ser necesaria en un sistema operativo cuando se necesite mantener una estructura lineal de datos pero no sea viable realizar las pausas significativas necesarias para redimensionar un arreglo dinámico. La variación en los tiempos individuales de ejecución al insertar elementos en un arreglo dinámico puede no ser aceptable, por ejemplo, dentro de una llamada al sistema que interrumpe al programa que la llama, porque los tiempos de pausa serían impredecibles y podrían causar problemas con tareas altamente interactivas, como la reproducción de audio.

En resumen, las listas enlazadas tienen la ventaja de permitir inserciones del orden de $\Theta{}(1)$ en posiciones arbitrarias conocidas. También ofrecen operaciones con tiempos de ejecución predecibles y consistentes. Pero necesitan usar más memoria, no permiten acceder a una posición arbitraria de la lista en tiempo constante y pueden ser significativamente más lentas debido a que cada elemento se almacena en una dirección de memoria totalmente distinta, desaprovechando el caché. La ventaja de los arreglos dinámicos es no tener ninguno de esos problemas, pero se sacrifica la previsibilidad y la velocidad de la inserción en posiciones arbitrarias.

\section*{Bibliografía}
\addcontentsline{toc}{section}{Bibliografía}

Cormen, T. H., Leiserson, C. E., Rivest, R. L., y Stein, C. (2022). \textit{Introduction to algorithms} (4.\textordfeminine{} ed.). MIT Press.
